\begin{abstract}

As part of a larger research agenda focused on mastery learning,
we were inspired by the benefits of Computer-Based Testing Facilities (CBTFs)
reported by
the University of Illinois at Urbana-Champaign,
and in Fall 2023 we decided to pilot such a facility ourselves.
If successful, we intended to scale up and eventually serve
the whole campus, not just one or two CS courses. 
Three years later, our CBTF operation has grown to 20x 
(geometric mean of 120\% growth in seated exam-hours per semester), and serves
departments across campus.
Bootstrapping the initial effort led to important lessons---some
evident only with distant hindsight and the need to scale---from which we hope others
might benefit.
Our two main takeaways are as follows.
First, even in a small pilot, avoid adopting
processes that won't scale,
%% Exceptional situations should create $O(1)$ work per
%%   student or per instructor rather than $O(n)$ work for CBTF staff,
since student and faculty expectations are heavily influenced by the
  policies and procedures established during the pilot phase.
Second, while various campus units typically ``own'' lab facilities,
disabled-student proctoring,
instructional technology, room reservations, and so on, usually
  no one ``owns'' whole-course proctored exam administration.  You will have to create new
  administrative processes, disrupt existing ones, correct mistaken
  faculty assumptions, and coordinate across nonporous administrative
  boundaries.

\end{abstract}

\begin{comment}
  %% submitted abstract plain text:


As part of a research program on mastery learning, inspired by the benefits of Computer-Based Testing reported by the University of Illinois at Urbana-Champaign, in Fall 2023 we ran a pilot of such a facility ourselves. If the pilot succeeded, we intended to scale up and eventually serve the whole campus, not just one or two CS courses. Three years later, our CBTF operation is 20x larger and serving departments across campus; bootstrapping the initial effort led to important lessons---some evident only with long hindsight---from which we hope others might benefit. Our two main takeaways are as follows. First, even in a small pilot, avoid adopting processes that won't scale, since student and faculty expectations are heavily influenced by the initial policies and procedures. Second, while various campus units typically ``own'' lab facilities, proctoring for students with disabilities, instructional software and hardware maintenance, room reservations, and so on, usually no one "owns" whole-course proctored exam administration.  You will have to create new administrative processes, disrupt existing ones, correct mistaken faculty assumptions, and coordinate across nonporous administrative boundaries.


\end{comment}
